\chapter{Threats to Validity}
\label{ch:threats}

\section{Internal Validity}

\paragraph{The policy influences the data it then learns from.} Within a
session the schedule decides which tests run, which decides what the store
contains, which decides the next schedule. Two things limit the damage.
Round~$1$ is identical across configurations under common random numbers, so
none starts from an advantage; and the per-test cost table is established by a
dedicated full-scan audit before any configuration runs, so no configuration's
ordering can contaminate the costs every configuration is scored against.

\paragraph{Instrument faults that produce numbers rather than errors.} Three
failure modes would yield plausible wrong values instead of a crash: a
(patch, test) pair executed and counted twice, a cached outcome disagreeing
with a fresh one, and a timeout entered into the cost model at the clock value.
Each is instrumented --- an accounting identity checked per scan, a halt on
disagreement, and quarantine of timeouts --- rather than asserted in prose.

\paragraph{Two tasks were dropped after the run.} The merge audit refused two
of the $99$ tasks that produced sessions, because the audit pass had not
recorded an outcome for every (patch, test) pair they needed. They are dropped
from every configuration rather than some, so the comparison stays paired ---
but the admitted set and the analysed set are not the same $100$ tasks.

\section{External Validity}

\paragraph{One benchmark, one language.} All results come from the Python
subset of ConDefects: single-file competitive-programming submissions with
whole-program, stdout-compared tests. Two consequences follow. Test
\emph{selection} is unavailable, so prioritisation is the only scheduling
lever, whereas in project-scale Java repair selection is available and has been
found to dominate~\citep{lou2024aprregression}. And contest pools are
deliberately heterogeneous in cost, which is favourable to the cost term --- an
interpreted language exaggerates the spread further. A uniform-cost corpus
would show less, which is condition C2 of Section~\ref{sec:disc-limits}.

\paragraph{One proposer model, at pilot scale.} Every candidate comes from
\code{qwen2.5-\hspace{0pt}coder:\hspace{0pt}7b} at one temperature, so the
rejection structure the store learns is one model's: which tests are
discriminating depends on which mistakes the model makes. The budget of twenty
rounds also truncates the learning curve, and
Section~\ref{sec:res-crossover} shows that it truncates it at the worst
possible place: \code{sched} only overtakes \code{frtp} at round $10$ and the
gap between them is still widening when the budget ends. A longer budget would
very likely show a larger effect, but the pilot cannot say how much larger,
and the direction of that bias should be read as working \emph{against} the
reported $-2.9\,\%$ rather than for it.

\paragraph{The $100$ tasks are a draw, not a random sample.} They were taken
from a corpus of roughly $1{,}244$ faults by the pool order fixed at freeze
time. Nothing in the analysis corrects for that, and the full-corpus run is the
only thing that removes it.

\section{Construct Validity}

\paragraph{The per-test cost table in the analysis is reconstructed.} The
report's corpus statistics (Table~\ref{tab:res-arena}) recover each test's
frozen cost by least squares from the replayed rounds, because a round's CPU is
the sum of the costs of the tests it ran. This is exact where a test appears in
enough rounds and noisy where it appears in few. It affects the descriptive
statistics and the cost-spread split of Table~\ref{tab:res-cv}; it does not
affect any head-to-head total, which is summed from the rounds themselves.

\paragraph{``Duplicate'' is defined by normalised text.} Two candidates that
differ in a variable name but express the same repair hypothesis are not
duplicates under this definition, although they are equally wasteful. The
duplicate rate is therefore a \emph{lower bound} on redundancy. The same caveat
applies to the $65\,\%$ figure reported for
\textsc{ChatRepair}~\citep{chen2025conversational}, which this work uses as its
reference point.

\paragraph{Cost is CPU; acceptance is plausibility.} A deployment paying for
wall clock on a loaded machine, or for a test service priced per invocation,
would prefer a different schedule from one paying for CPU --- and the units can
disagree by more than the effect being measured
(Section~\ref{sec:disc-units}), which is why both are reported for every
configuration. Separately, a candidate surviving the whole pool is
\emph{accepted}, and the pool is finite: passing a finite suite is weaker than
being correct~\citep{qi2015plausibility}. That does not affect the cost claims,
which hold the acceptance criterion fixed across configurations, but it bounds
what the repair-rate numbers mean.

\paragraph{The sandbox clock shapes the corpus silently.} A wall-clock kill is
applied to the reference program first, so a test the correct program cannot
finish leaves the pool, and a task whose pool empties leaves the corpus. This
keeps the oracle sound, but it removes faults non-randomly, from the expensive
end --- the end where the cost term would matter most.

\section{Conclusion Validity}

\paragraph{Sample size and skew.} Ninety-seven tasks with three seeds is small
for an effect whose per-task distribution is heavily right-skewed and dominated
by a few expensive tasks: the top third by cost spread carries $79\,\%$ of the
corpus's validation CPU. Comparisons therefore use the Wilcoxon signed-rank
test on per-task totals rather than a mean-based test, every point estimate is
reported with an interval, and the contrast where the aggregate and the
per-task counts disagree is reported both ways
(Table~\ref{tab:res-frtp-split}).

\paragraph{Clustering.} The three seeds of one task share the task, the pool
and the cost structure, so treating them as independent would understate every
interval by roughly $\sqrt{3}$. Intervals come from a cluster bootstrap that
resamples whole tasks.

\paragraph{Multiple comparisons.} Four contrasts are declared in advance and
carry a Holm correction. The ablation ladder, the cost-spread split and the
coverage-check replay are reported with intervals and labelled exploratory, and
no conclusion rests on one of them alone.

\paragraph{A significant $p$ and an interval covering zero are not a
contradiction.} For \code{sched} against \code{frtp} the Wilcoxon test is
significant ($p = 0.026$ after correction) while the bootstrap interval on
total CPU covers zero. They measure different things: the test asks which
configuration is cheaper on a \emph{typical} task, and the answer is \frtp, on
$45$ tasks against $18$; the interval asks about the \emph{total}, which
\code{sched} wins because its few wins are large. Both are reported, because
reporting either alone would be misleading.
